\section{Peta Konsep Pemrograman Game}

Mata kuliah ini mencakup 15 bab yang saling terkait:

\begin{enumerate}
  \item \textbf{Bab II}: Landasan Teori dan Konsep Dasar -- MDA, genre, GDD, lifecycle
  \item \textbf{Bab III}: Pengenalan Unity Engine -- Instalasi, Editor, workflow
  \item \textbf{Bab IV}: Dasar Pemrograman C\# untuk Unity -- MonoBehaviour, lifecycle, C\#
  \item \textbf{Bab V}: GameObject, Component, Prefab, Scene -- Fondasi Unity
  \item \textbf{Bab VI}: Input System dan Player Controller -- Movement, camera
  \item \textbf{Bab VII}: Physics, Collision, dan Rigidbody -- Sistem fisika
  \item \textbf{Bab VIII}: Animation dan Animator -- Animasi karakter
  \item \textbf{Bab IX}: UI/UX Game -- Canvas, HUD, menu
  \item \textbf{Bab X}: Audio dalam Game -- AudioSource, BGM, SFX
  \item \textbf{Bab XI}: Level Design dan Game Mechanics -- Desain level, mekanik
  \item \textbf{Bab XII}: AI Sederhana untuk Game -- Enemy AI, pathfinding
  \item \textbf{Bab XIII}: Optimization dan Debugging -- Profiler, best practices
  \item \textbf{Bab XIV}: Build dan Deployment -- Build settings, platform
  \item \textbf{Bab XV}: Evaluasi, Refleksi, dan Integrasi Kompetensi
\end{enumerate}

\textbf{Alur Pembelajaran:}
\begin{itemize}
  \item Bab II: Fondasi teoretis game development
  \item Bab III--V: Pengenalan Unity dan struktur dasar
  \item Bab VI--X: Implementasi fitur game (input, physics, animasi, UI, audio)
  \item Bab XI--XII: Level design dan AI
  \item Bab XIII--XIV: Optimasi dan deployment
  \item Bab XV: Evaluasi dan integrasi
\end{itemize}
